\chapter{Tailwind CSS for MERN Apps}

\begin{notebox}[NOTE]
This chapter does not teach CSS fundamentals (the box model, flexbox theory,
specificity, etc.). It assumes you already know CSS and focuses only on how
to apply Tailwind's utility classes to build the screens of the Task
Management System quickly and consistently.
\end{notebox}

\section{Responsive Prefixes}

Tailwind applies styles at a breakpoint using a prefix: \texttt{sm:}
(640px+), \texttt{md:} (768px+), \texttt{lg:} (1024px+). Unprefixed classes
apply at all sizes; a prefixed class overrides it from that breakpoint up.

\begin{lstlisting}
<div className="flex flex-col md:flex-row gap-4">
  {/* stacked on mobile, side-by-side from tablet up */}
</div>
\end{lstlisting}

\section{Flex, Grid, and Spacing}

\begin{lstlisting}
{/* Flex: horizontal toolbar */}
<div className="flex items-center justify-between p-4">
  <h1 className="text-xl font-bold">Tasks</h1>
  <button className="px-4 py-2 bg-blue-600 text-white rounded">New Task</button>
</div>

{/* Grid: responsive task card grid */}
<div className="grid grid-cols-1 sm:grid-cols-2 lg:grid-cols-3 gap-4">
  {tasks.map((t) => <TaskCard key={t._id} task={t} />)}
</div>
\end{lstlisting}

Common spacing utilities: \texttt{p-4} (padding), \texttt{m-4} (margin),
\texttt{gap-4} (flex/grid gap), \texttt{space-y-4} (vertical gap between
stacked children without a flex/grid parent).

\section{Styling Forms}

\begin{lstlisting}
<div>
  <label className="block text-sm font-medium text-gray-700 mb-1">Email</label>
  <input
    type="email"
    className="w-full px-3 py-2 border border-gray-300 rounded-md
               focus:outline-none focus:ring-2 focus:ring-blue-500
               focus:border-transparent"
  />
</div>
\end{lstlisting}

\texttt{focus:ring-2} and \texttt{focus:ring-blue-500} give a visible
focus ring without relying on the browser's default outline --- useful for
consistent branding, but keep it visible for keyboard accessibility.

\section{Buttons: Primary / Secondary / Danger}

Compose a base class string and merge variant-specific classes, so every
button shares padding/rounding/transition but differs only in color.

\begin{lstlisting}[language=JavaScript]
// components/Button.jsx
const base = "px-4 py-2 rounded-md font-medium transition-colors disabled:opacity-50";

const variants = {
  primary: "bg-blue-600 text-white hover:bg-blue-700",
  secondary: "bg-gray-100 text-gray-800 hover:bg-gray-200",
  danger: "bg-red-600 text-white hover:bg-red-700",
};

function Button({ variant = "primary", children, ...props }) {
  return (
    <button className={`${base} ${variants[variant]}`} {...props}>
      {children}
    </button>
  );
}

export default Button;

// <Button variant="primary">Save</Button>
// <Button variant="secondary">Cancel</Button>
// <Button variant="danger">Delete</Button>
\end{lstlisting}

\section{Cards: TaskCard}

\begin{lstlisting}[language=JavaScript]
// components/TaskCard.jsx
function TaskCard({ task }) {
  const statusColor = {
    pending: "bg-yellow-100 text-yellow-800",
    "in-progress": "bg-blue-100 text-blue-800",
    done: "bg-green-100 text-green-800",
  };

  return (
    <div className="border border-gray-200 rounded-lg p-4 shadow-sm
                     hover:shadow-md transition-shadow bg-white">
      <div className="flex justify-between items-start mb-2">
        <h3 className="font-semibold text-gray-900">{task.title}</h3>
        <span className={`text-xs px-2 py-1 rounded-full ${statusColor[task.status]}`}>
          {task.status}
        </span>
      </div>
      <p className="text-sm text-gray-600 line-clamp-2">{task.description}</p>
    </div>
  );
}

export default TaskCard;
\end{lstlisting}

\section{Modals}

A modal overlay uses a fixed, full-screen backdrop with the dialog centered
on top of it.

\begin{lstlisting}[language=JavaScript]
// components/Modal.jsx
function Modal({ isOpen, onClose, children }) {
  if (!isOpen) return null;

  return (
    <div className="fixed inset-0 bg-black/50 flex items-center justify-center z-50">
      <div className="bg-white rounded-lg shadow-xl w-full max-w-md p-6 relative">
        <button
          onClick={onClose}
          className="absolute top-3 right-3 text-gray-400 hover:text-gray-600"
        >
          &times;
        </button>
        {children}
      </div>
    </div>
  );
}

export default Modal;

// <Modal isOpen={showModal} onClose={() => setShowModal(false)}>
//   <TaskForm onSaved={() => setShowModal(false)} />
// </Modal>
\end{lstlisting}

\begin{notebox}[NOTE]
\texttt{fixed inset-0} pins the overlay to all four edges of the viewport
regardless of scroll position. \texttt{bg-black/50} is Tailwind's slash
opacity syntax for black at 50\% opacity --- no separate opacity class
needed.
\end{notebox}

\section{Navbar with Responsive Menu}

\begin{lstlisting}[language=JavaScript]
// components/Navbar.jsx
import { useState } from "react";

function Navbar() {
  const [open, setOpen] = useState(false);

  return (
    <nav className="bg-white border-b border-gray-200 px-4 py-3">
      <div className="flex items-center justify-between">
        <span className="font-bold text-lg text-blue-600">TaskManager</span>

        {/* Desktop links */}
        <div className="hidden md:flex items-center gap-6">
          <a href="/tasks" className="text-gray-700 hover:text-blue-600">Tasks</a>
          <a href="/profile" className="text-gray-700 hover:text-blue-600">Profile</a>
        </div>

        {/* Mobile toggle */}
        <button className="md:hidden" onClick={() => setOpen(!open)}>
          Menu
        </button>
      </div>

      {/* Mobile links */}
      {open && (
        <div className="md:hidden mt-3 flex flex-col gap-2">
          <a href="/tasks" className="text-gray-700">Tasks</a>
          <a href="/profile" className="text-gray-700">Profile</a>
        </div>
      )}
    </nav>
  );
}

export default Navbar;
\end{lstlisting}

\section{Sidebar Layout}

\begin{lstlisting}
<div className="flex min-h-screen">
  <aside className="hidden lg:block w-64 bg-gray-900 text-white p-4">
    {/* nav links */}
  </aside>
  <main className="flex-1 p-6 bg-gray-50">
    {/* page content */}
  </main>
</div>
\end{lstlisting}

\texttt{hidden lg:block} hides the sidebar below the \texttt{lg} breakpoint
and shows it from 1024px up, so mobile users get the toolbar/nav instead of
a permanent sidebar eating screen space.

\begin{tipbox}[TIP]
Combine \texttt{hidden md:flex} (or \texttt{lg:block}) with a mobile toggle
button anywhere you need "desktop shows X, mobile shows a menu button
instead" --- it's the single most reused responsive pattern across
navbars, sidebars, and filter panels.
\end{tipbox}

\whatwhyhow
{Tailwind utility classes (layout, spacing, color, state variants like \texttt{focus:} and \texttt{hover:}, and responsive prefixes like \texttt{md:}) are composed directly in JSX to build buttons, cards, modals, navbars, and responsive layouts.}
{Utility classes let you style components without maintaining separate CSS files or fighting specificity, and responsive prefixes make one component adapt across breakpoints without media query boilerplate.}
{Build small reusable components (\texttt{Button}, \texttt{TaskCard}, \texttt{Modal}) with variant classNames, and combine \texttt{hidden}/\texttt{flex}/\texttt{grid} with breakpoint prefixes for layouts that adapt from mobile to desktop.}
